% ===================================================================
%  ਸਾਰਣੀ ਨੰ. 16 — ਵੱਡੀ ਦੁਕਾਨ। — ਮੰਡੀ। — ਖਿਡੌਣੇ।
%  Traduction 2026 d'apres texto/io, controlee sur texto/fr, texto/en
%  et texto/hi. Consigne : voir texto/pa/10-sarni-01.tex.
%
%  Le tableau d'astronomie (bloc 15-1) porte des renvois a LETTRES,
%  (a) a (f), graves sur la planche elle-meme : ils se rangent sous
%  le numero 85, comme le dit gravuri/literi.json. L'ido y saute la
%  parenthese ouvrante du (d) ; la traduction la donne, et le
%  controle passe ce bloc sans rien dire.
% ===================================================================

%%K t16-apar-1 apar
\VUsaut{4.0mm}
\VUtitre{15.0pt}{60}{ਸਾਰਣੀ ਨੰ. 16}
\VUsaut{2.0mm}
\VUpk{13.2pt}{40}{ਵੱਡੀ ਦੁਕਾਨ। --- ਮੰਡੀ।}
\VUpk{13.2pt}{40}{ਖਿਡੌਣੇ।}
\VUsaut{3.4mm}

%%K t16-01-1 p
1. --- ਨਵਾਂ ਵਰ੍ਹਾ ਆ ਰਿਹਾ ਹੈ। ਵੱਡੀਆਂ ਦੁਕਾਨਾਂ ਨੇ \VUgras{ਖਿਡੌਣਿਆਂ}
\textsuperscript{(1)} ਤੇ ਨਵੇਂ ਵਰ੍ਹੇ ਦੀਆਂ ਸੌਗਾਤਾਂ ਦੀ ਆਪਣੀ ਸਜਾਵਟ ਤਿਆਰ ਕਰ ਲਈ ਹੈ।

%%K t16-01-2 p
ਮੈਂ ਆਪਣੇ ਦਾਦਾ ਨੂੰ ਆਖਿਆ ਕਿ ਮੈਨੂੰ ਇਹ ਸੋਹਣੀ ਸਜਾਵਟ ਵੇਖਣ ਮੰਡੀ ਲੈ ਚੱਲੋ, ਤੇ ਉਹ ਮੰਨ ਗਏ।

%%K t16-02-1 p
2. --- ਪਹਿਲਾਂ ਅਸੀਂ ਹਥਿਆਰਾਂ ਦੀਆਂ ਕੁਝ ਸੋਹਣੀਆਂ ਸਜਾਵਟਾਂ ਦੇ ਸਾਹਮਣੇ ਰੁਕੇ, ਜਿਨ੍ਹਾਂ ਵਿਚ ਇਕ
\VUgras{ਰਾਈਫ਼ਲ}, ਇਕ \VUgras{ਝਾਲਰ ਵਾਲੀ ਟੋਪੀ}, ਇਕ \VUgras{ਕਿਰਚ}, ਇਕ
\VUgras{ਤਲਵਾਰ ਦੀ ਪੇਟੀ} ਤੇ ਇਕ ਜੋੜੀ \VUgras{ਮੋਢੇ ਦੇ ਫੁੰਮਣ} ਸਨ, ਜਾਂ ਫਿਰ ਇਕ
\VUgras{ਖੋਦ}, ਇਕ \VUgras{ਛਾਤੀ ਦਾ ਕਵਚ} \textsuperscript{(3)} ਤੇ ਇਕ
\VUgras{ਪਿਸਤੌਲ} \textsuperscript{(4)}। ਇਹ ਸਭ ਮੈਨੂੰ \VUgras{ਜਾਦੂਈ ਲਾਲਟੈਣਾਂ}
\textsuperscript{(2)} ਨਾਲੋਂ ਕਿਤੇ ਵੱਧ ਭਾਇਆ।

%%K t16-tit-1 sub
\VUsaut{3.0mm}
\VUpk{10.2pt}{40}{ਵੇਖਣ ਜੋਗ ਚੀਜ਼ਾਂ।}
\VUsaut{2.6mm}

%%K t16-03-1 p
3. --- ਉਸੇ ਸ਼ੋਅਬੇ ਵਿਚ ਇਕ \VUgras{ਪਹਿਰੇਦਾਰ} \textsuperscript{(5)} ਆਪਣੀ
\VUgras{ਪਹਿਰੇ ਦੀ ਕੋਠੜੀ} ਵਿਚ ਸੀ; ਉਹ ਗੱਤੇ ਦੇ ਪੂਰੇ ਚਿੜੀਆਘਰ ਉੱਤੇ ਪਹਿਰਾ ਦਿੰਦਾ ਲਗਦਾ ਸੀ।
ਕਿੰਨੇ ਜਾਨਵਰ ਸਨ ਉੱਥੇ! ਆਪਣੇ ਅੱਡੇ ਉੱਤੇ ਇਕ \VUgras{ਤੋਤਾ} \textsuperscript{(6)}, ਨੱਕ
ਉੱਤੇ ਸਿੰਗ ਵਾਲਾ ਇਕ \VUgras{ਗੈਂਡਾ} \textsuperscript{(7)}, ਇਕ \VUgras{ਬਾਰਾਂਸਿੰਗਾ}
\textsuperscript{(8)}, ਇਕ \VUgras{ਸ਼ੁਤਰਮੁਰਗ਼} \textsuperscript{(9)}, ਇਕ
\VUgras{ਬਾਂਦਰ} \textsuperscript{(10)} ਜੋ ਅਖ਼ਰੋਟ ਤੋੜ ਰਿਹਾ ਸੀ, ਇਕ \VUgras{ਲੂੰਬੜੀ}
\textsuperscript{(11)} ਜੋ ਇਕ ਕੁਕੜੀ ਚੁੱਕੀ ਲੈ ਜਾਂਦੀ ਸੀ, ਇਕ \VUgras{ਮਗਰਮੱਛ}
\textsuperscript{(12)} ਆਪਣੇ ਵੱਡੇ ਜਬਾੜੇ ਖੋਲ੍ਹੀ, ਇਕ \VUgras{ਊਠ}
\textsuperscript{(13)}, ਇਕ ਸਜੀਲਾ \VUgras{ਗਰੇਹਾਊਂਡ} \textsuperscript{(14)}, ਇਕ
\VUgras{ਤੇਂਦੂਆ} \textsuperscript{(15)}, ਤੇ ਫਿਰ ਦੋ ਜਾਨਵਰ ਜਿਨ੍ਹਾਂ ਨੂੰ ਮੈਂ ਜਾਣਦਾ ਨਾ ਸੀ,
ਜੋ, ਮੈਨੂੰ ਦੱਸਿਆ ਗਿਆ, ਇਕ \VUgras{ਨਿਓਲਾ} \textsuperscript{(16)} ਤੇ ਇਕ
\VUgras{ਲੱਕੜਬੱਗਾ} \textsuperscript{(17)} ਹਨ। ਉੱਥੇ ਇਕ \VUgras{ਚੀਤਾ}
\textsuperscript{(18)} ਤੇ ਇਕ \VUgras{ਬਘਿਆੜ} \textsuperscript{(21)} ਵੀ ਸੀ; ਕੋਲ ਹੀ
ਪਿਆਰੇ ਨਿੱਕੇ \VUgras{ਚੂਹੇ} \textsuperscript{(19)} ਤੇ ਰਬੜ ਦੀਆਂ ਨਿੱਕੀਆਂ
\VUgras{ਚੂਹੀਆਂ} \textsuperscript{(20)} ਸਨ। ਅੰਤ ਵਿਚ ਇਕ \VUgras{ਦਰਿਆਈ ਘੋੜਾ}
\textsuperscript{(22)} ਤੇ ਇਕ ਕੁੱਬ ਵਾਲਾ \VUgras{ਊਠ} \textsuperscript{(23)} ਸੰਗ੍ਰਹਿ
ਪੂਰਾ ਕਰਦੇ ਸਨ। ਮੈਂ ਉੱਥੇ ਘੰਟਿਆਂ ਖੜਾ ਰਹਿੰਦਾ, ਜੇ ਕੋਲ ਹੀ \VUgras{ਟੀਨ ਦੇ ਸਿਪਾਹੀਆਂ}
\textsuperscript{(29)} ਨਾਲ ਭਰਿਆ ਇਕ ਸ਼ਾਨਦਾਰ ਪੜਾਅ ਮੇਰੀ ਅੱਖ ਨਾ ਖਿੱਚ ਲੈਂਦਾ।

%%K t16-tit-2 sub
\VUsaut{4.0mm}
\VUpk{10.2pt}{40}{ਸਿਪਾਹੀ ਤੇ ਜੰਗ।}
\VUsaut{2.8mm}

%%K t16-04-1 p
4. --- ਮੇਰੇ ਦਾਦਾ, ਜੋ \VUgras{ਅਪਾਹਜ ਸਾਬਕਾ ਸਿਪਾਹੀ} \textsuperscript{(24)} ਹਨ ਤੇ
ਜਿਨ੍ਹਾਂ ਨੂੰ ਇਕ \VUgras{ਜ਼ਖ਼ਮ} ਪਿੱਛੋਂ ਫ਼ੌਜ ਛੱਡਣੀ ਪਈ, ਜਿਸ ਨੇ ਉਨ੍ਹਾਂ ਨੂੰ
\VUgras{ਫ਼ੌਜੀ ਤਗ਼ਮਾ} ਦਿਵਾਇਆ, ਉਨ੍ਹਾਂ ਨੂੰ ਮੈਨੂੰ ਉਨ੍ਹਾਂ \VUgras{ਮੁਹਿੰਮਾਂ} ਦੀਆਂ
ਗੱਲਾਂ ਸੁਣਾਉਣਾ ਬਹੁਤ ਚੰਗਾ ਲਗਦਾ ਹੈ ਜਿਨ੍ਹਾਂ ਵਿਚ ਉਹ ਸ਼ਾਮਲ ਹੋਏ। ਉਹ ਮੈਨੂੰ ਉਸ ਨਿੱਕੀ ਫ਼ੌਜ
ਦੀਆਂ ਵੱਖੋ-ਵੱਖ ਚਾਲਾਂ ਸਮਝਾ ਕੇ ਬਹੁਤ ਖ਼ੁਸ਼ ਹੋਏ। ਮੰਡੀ ਵੇਖਣ ਆਏ ਅਸਲੀ ਸਿਪਾਹੀਆਂ ਦਾ ਇਕ ਟੋਲਾ —
ਇਕ \VUgras{ਇੰਜੀਨੀਅਰ ਸਿਪਾਹੀ} \textsuperscript{(25)}, ਇਕ \VUgras{ਪਹਾੜੀ ਹੌਲਾ ਸਿਪਾਹੀ}
\textsuperscript{(26)} ਆਪਣੀ \VUgras{ਟੋਪੀ} ਵਿਚ, ਪੈਦਲ ਫ਼ੌਜ ਦਾ ਇਕ
\VUgras{ਮੁੱਖ ਹੌਲਦਾਰ} \textsuperscript{(27)} ਤੇ ਤੋਪਖ਼ਾਨੇ ਦਾ ਇਕ \VUgras{ਹੌਲਦਾਰ}
\textsuperscript{(28)}, ਜਿਸ ਦੀ ਆਸਤੀਨ ਉੱਤੇ ਸੋਨੇ ਦੀ \VUgras{ਧਾਰੀ} ਸੀ — ਉਹ ਵੀ ਵੇਰਵਾ
ਸੁਣਨ ਸਾਡੇ ਕੋਲ ਰੁਕ ਗਏ।

%%K t16-05-1 p
5. --- ਮੇਰੇ ਦਾਦਾ, ਇੰਨੇ ਮਸ਼ਹੂਰ ਸੁਣਨ ਵਾਲੇ ਲੱਭ ਕੇ ਬਹੁਤ ਮਾਣ ਨਾਲ ਭਰੇ, ਉਨ੍ਹਾਂ ਝੱਟ ਜੰਗ ਦੀ
ਇਕ ਪੂਰੀ ਵਿਉਂਤ ਘੜ ਸੁੱਟੀ।

%%K t16-05-2 p
ਉਹ ਕਿਲ੍ਹਾਬੰਦ ਸ਼ਹਿਰ, ਆਪਣੀਆਂ \VUgras{ਫ਼ਸੀਲਾਂ} ਤੇ \VUgras{ਖਾਈਆਂ} ਸਮੇਤ,
\VUgras{ਵੈਰੀ ਫ਼ੌਜ} ਨਾਲ ਘਿਰਿਆ ਸੀ, ਜੋ ਲੰਮੀ \VUgras{ਗੋਲਾਬਾਰੀ} ਪਿੱਛੋਂ ਉਸ ਉੱਤੇ
\VUgras{ਧਾਵਾ} ਬੋਲਣ ਦੀ ਤਿਆਰੀ ਵਿਚ ਸੀ। ਪਰ \VUgras{ਘਿਰੇ ਹੋਏ} ਲੋਕ, ਭੁੱਖ ਤੋਂ ਮਜਬੂਰ,
\VUgras{ਘੇਰਨ ਵਾਲਿਆਂ} ਦੇ \VUgras{ਪੜਾਅ} ਉੱਤੇ ਇਕ \VUgras{ਬਾਹਰ-ਵਾਰ} ਕਰ ਬੈਠੇ ਸਨ।
\VUgras{ਤੰਬੂਆਂ} ਦੇ ਚੁਫੇਰੇ ਇਕ ਤਿੱਖੀ \VUgras{ਲੜਾਈ} ਵਿਚ \VUgras{ਹੱਲੇ} ਨੂੰ ਪਿੱਛੇ ਧੱਕ
ਕੇ \VUgras{ਜਿੱਤਣ ਵਾਲੇ} ਘੇਰਨ ਵਾਲਿਆਂ ਨੇ ਆਪਣੇ \VUgras{ਘੋੜਸਵਾਰ ਦਸਤੇ}
\VUgras{ਹਾਰੇ ਹੋਏ} ਹਮਲਾਵਰਾਂ ਦਾ \VUgras{ਪਿੱਛਾ} ਕਰਨ ਘੱਲੇ। ਉਹ \VUgras{ਪਿੱਛੇ ਹਟੇ},
\VUgras{ਰਣ-ਭੋਇੰ} ਨੂੰ \VUgras{ਮੁਰਦਿਆਂ} ਤੇ \VUgras{ਜ਼ਖ਼ਮੀਆਂ} ਨਾਲ ਵਿਛਾ ਕੇ।
\VUgras{ਸਟਰੈਚਰ ਵਾਲੇ} ਜ਼ਖ਼ਮੀਆਂ ਨੂੰ \VUgras{ਚਲਦੇ ਦਵਾਖ਼ਾਨੇ} ਤਕ ਲੈ ਗਏ, ਤਾਂ ਜੋ
\VUgras{ਜਰਾਹ} ਉਨ੍ਹਾਂ ਦੀ ਸੰਭਾਲ ਕਰਨ।

%%K t16-05-3 p
ਕੁਝ \VUgras{ਬਟਾਲੀਅਨਾਂ} ਦੇ ਸੂਰਮੇ ਮੁਕਾਬਲੇ ਦੇ ਬਾਵਜੂਦ, ਜਿਨ੍ਹਾਂ ਨੇ \VUgras{ਚੌਕੋਰ ਵਿਊਹ}
ਬਣਾਇਆ ਸੀ, \VUgras{ਹਾਰ} ਪੂਰੀ ਸੀ; ਬਚੀ ਫ਼ੌਜ \VUgras{ਨੱਸ ਗਈ}, ਬਹੁਤ ਸਾਰੇ
\VUgras{ਬੰਦੀ} ਪਿੱਛੇ ਛੱਡ ਕੇ। ਵੈਰੀ ਸ਼ਹਿਰ ਲੈ ਕੇ ਆਪਣੀ \VUgras{ਜਿੱਤ} ਉੱਤੇ ਤਾਜ ਰੱਖਣ ਨੂੰ
ਸੀ। ਸ਼ਾਇਦ ਇਹੋ ਮੁਹਿੰਮ ਦਾ ਅੰਤ ਹੁੰਦਾ, ਤੇ ਇਕ \VUgras{ਜੰਗਬੰਦੀ} ਪਿੱਛੋਂ ਇਕ
\VUgras{ਅਮਨ ਦੇ ਸਮਝੌਤੇ} ਉੱਤੇ ਦਸਤਖ਼ਤ ਹੋ ਸਕਦੇ।

%%K t16-tit-3 sub
\VUsaut{4.4mm}
\VUpk{10.2pt}{40}{ਨਵੇਂ ਵਰ੍ਹੇ ਦੀਆਂ ਸੌਗਾਤਾਂ।}
\VUsaut{2.8mm}

%%K t16-06-1 p
6. --- ਉਹ ਜਵਾਨ ਸਿਪਾਹੀ, ਜੋ ਸਾਬਕਾ ਸਿਪਾਹੀ ਦੀ ਰੰਗੀਲੀ ਕਹਾਣੀ ਸੁਣਦੇ ਹੋਏ ਹੱਸ ਰਹੇ ਸਨ,
ਉਨ੍ਹਾਂ ਨੇ ਉਨ੍ਹਾਂ ਕੋਲੋਂ ਕੁਝ ਹੋਰ ਸਵਾਲ ਪੁੱਛੇ। ਰਿਹਾ ਮੈਂ, ਤਾਂ ਮੈਂ ਕਾਊਂਟਰ ਦਾ ਚੱਕਰ ਕੱਟ ਕੇ
ਇਕ ਸੋਹਣੀ \VUgras{ਕਮਾਨ} \textsuperscript{(32)} ਤੇ ਇਕ \VUgras{ਧਨੁਸ਼}
\textsuperscript{(33)} ਉਸ ਦੇ \VUgras{ਤੀਰ} \textsuperscript{(31)} ਸਮੇਤ ਵੇਖਣ ਗਿਆ। ਉੱਥੇ
ਮੈਨੂੰ ਜਾਨਵਰਾਂ ਦਾ ਇਕ ਹੋਰ ਸੰਗ੍ਰਹਿ ਮਿਲਿਆ: ਦੋ \VUgras{ਗਾਉਣ ਵਾਲੀਆਂ ਚਿੜੀਆਂ}
\textsuperscript{(34)}, ਇਕ ਚਾਬੀ ਵਾਲੀ \VUgras{ਬੁਲਬੁਲ} ਤੇ \VUgras{ਫੁਦਕੀ} ਜੋ ਗਲਾ
ਪਾੜ ਕੇ ਗਾ ਰਹੀਆਂ ਸਨ, ਤੇ ਅਨੋਖੇ ਜੀਵਾਂ ਦੀ ਇਕ ਕਤਾਰ — ਇਕ \VUgras{ਚਮਗਿੱਦੜ}
\textsuperscript{(35)}, ਇਕ \VUgras{ਅਜਗਰ} \textsuperscript{(36)}, ਇਕ \VUgras{ਕੱਛੂ}
\textsuperscript{(37)}, ਇਕ \VUgras{ਸੀਲ} \textsuperscript{(38)}, ਇਕ \VUgras{ਵ੍ਹੇਲ}
\textsuperscript{(39)} ਤੇ ਸੋਨੇ ਦੀ ਝਿੱਲੀ ਨਾਲ ਬਣੀ ਇਕ \VUgras{ਸ਼ਾਰਕ}
\textsuperscript{(40)}।

%%K t16-07-1 p
7. --- ਮੈਂ ਉਨ੍ਹਾਂ ਨੂੰ ਵੇਖਣ ਵੱਧ ਦੇਰ ਨਾ ਰੁਕਿਆ, ਕਿਉਂਕਿ ਮੇਰੀ ਅੱਖ ਉਸ ਚੀਜ਼ ਉੱਤੇ ਪੈ ਚੁੱਕੀ
ਸੀ ਜਿਸ ਦਾ ਮੈਂ ਸੁਪਨਾ ਵੇਖਦਾ ਸੀ: ਇਕ ਸ਼ਾਨਦਾਰ \VUgras{ਕਠਪੁਤਲੀ ਰੰਗ-ਮੰਚ}
\textsuperscript{(41)}, ਸ਼ਾਨਦਾਰ \VUgras{ਦ੍ਰਿਸ਼-ਪਟ} ਤੇ ਇਕ ਸੋਹਣੇ ਲਾਲ
\VUgras{ਪਰਦੇ} ਸਮੇਤ। \VUgras{ਮੰਚ} ਉੱਤੇ ਦੋ ਅਦਾਕਾਰ ਕਿਸੇ ਬਹੁਤ ਦਿਲਚਸਪ
\VUgras{ਨਾਟਕ} ਵਿਚ ਆਪਣੀ \VUgras{ਭੂਮਿਕਾ} ਨਿਭਾਉਂਦੇ ਲਗਦੇ ਸਨ, ਕਿਉਂਕਿ ਡੱਬਿਆਂ ਦੇ ਗੱਤੇ
ਦੇ ਨਿੱਕੇ \VUgras{ਵੇਖਣ ਵਾਲੇ}, ਜਾਂ ਜੋ \VUgras{ਅਗਲੀਆਂ ਕੁਰਸੀਆਂ} ਤੇ
\VUgras{ਪਿਛਲੇ ਹਿੱਸੇ} ਵਿਚ \VUgras{ਸਾਜ਼-ਚਾਲਕ} ਦੇ ਪਿੱਛੇ ਬੈਠੇ ਸਨ, ਉਹ ਗਰਮਜੋਸ਼ੀ ਨਾਲ
ਤਾੜੀ ਵਜਾ ਰਹੇ ਸਨ।

%%K t16-07-2 p
ਅਫ਼ਸੋਸ, ਉਹ ਰੰਗ-ਮੰਚ ਬਹੁਤ ਮਹਿੰਗਾ ਸੀ।

%%K t16-08-1 p
8. --- ਇਸ ਤੋਂ ਬਿਨਾਂ, ਖਿਡੌਣਿਆਂ ਵਿਚੋਂ ਚੁਣਨਾ ਔਖਾ ਸੀ, ਕਿਉਂਕਿ ਖਿੜਕੀਆਂ ਹਰ ਤਰ੍ਹਾਂ ਦੇ
ਖਿਡੌਣਿਆਂ ਨਾਲ ਭਰੀਆਂ ਸਨ: \VUgras{ਹਾਰਲੇਕਿਨ} \textsuperscript{(42)},
\VUgras{ਪੋਲੀਸ਼ੀਨੇਲ} \textsuperscript{(43)}, \VUgras{ਪਿਏਰੋ} \textsuperscript{(44)},
ਖੰਭ ਫੜਫੜਾਉਂਦੇ \VUgras{ਉੱਲੂ} \textsuperscript{(45)}, ਸੀਟੀ ਵਜਾਉਂਦੇ
\VUgras{ਕਾਲਕੰਠ} \textsuperscript{(46)}, ਭੌਂਕਦੇ \VUgras{ਪੂਡਲ}, ਇਕ
\VUgras{ਗ੍ਰਾਮੋਫ਼ੋਨ} \textsuperscript{(47)} ਤੇ \VUgras{ਬੁਝਾਰਤਾਂ}। ਇਨ੍ਹਾਂ ਵਿਚੋਂ ਇਕ
ਖਿਡੌਣਾ ਮੈਨੂੰ ਬਹੁਤ ਭਾਇਆ: ਉਸ ਵਿਚ ਇਕ \VUgras{ਸਮੁੰਦਰੀ ਲੜਾਈ} \textsuperscript{(48)}
ਵਿਖਾਈ ਗਈ ਸੀ, ਜਿਸ ਵਿਚ \VUgras{ਨਾਰੀਅਲ ਦੇ ਰੁੱਖਾਂ} ਨਾਲ ਭਰੇ ਕੰਢੇ ਦੇ ਸਾਹਮਣੇ
\VUgras{ਸਮੁੰਦਰੀ ਡਾਕੂ} ਇਕ \VUgras{ਜਹਾਜ਼} ਉੱਤੇ ਹੱਲਾ ਕਰਦੇ ਹਨ।

%%K t16-noto-1 noto
\VUnotes{143.2mm}{%
\textit{(*) ਸਾਰਣੀ ਨੰ. 6 ਵੇਖੋ।}%
}

%%K t16-09-1 p
9. --- ਮੈਂ ਇਹ ਸਾਰੀਆਂ ਹੈਰਾਨ ਕਰ ਦੇਣ ਵਾਲੀਆਂ ਚੀਜ਼ਾਂ ਬੜੇ ਸਕੂਨ ਨਾਲ ਵੇਖ ਸਕਿਆ, ਕਿਉਂਕਿ
ਦੁਕਾਨ ਵਿਚ ਲੋਕ ਘੱਟ ਹੀ ਸਨ — ਮੁਸ਼ਕਲ ਨਾਲ ਮੁੱਠ ਭਰ ਤਮਾਸ਼ਬੀਨ, ਇਕ
\VUgras{ਘੋੜਸਵਾਰ ਸਿਪਾਹੀ} \textsuperscript{(49)}, ਤੇ ਇਕ \VUgras{ਵਸੂਲੀ ਕਰਨ ਵਾਲਾ}
\textsuperscript{(50)}, ਜੋ ਮੁੱਖ \VUgras{ਗੱਲੇ} \textsuperscript{(51)} ਉੱਤੇ ਇਕ
\VUgras{ਹੁੰਡੀ} ਪੇਸ਼ ਕਰ ਰਿਹਾ ਸੀ।

%%K t16-10-1 p
10. --- ਮੈਂ ਆਪਣੇ ਦਾਦਾ ਕੋਲੋਂ ਪੁੱਛਿਆ ਕਿ \VUgras{ਖ਼ਜ਼ਾਨਚੀ} ਦੇ ਪਿੱਛੇ ਤਾਕਾਂ ਉੱਤੇ ਰੱਖੀਆਂ
ਉਹ ਕਿਤਾਬਾਂ ਕਿਸ ਕੰਮ ਦੀਆਂ ਹਨ; ਉਨ੍ਹਾਂ ਦੱਸਿਆ ਕਿ ਉਹ \VUgras{ਹਿਸਾਬ ਦੀਆਂ ਵਹੀਆਂ} ਹਨ।

%%K t16-tit-4 sub
\VUsaut{3.6mm}
\VUpk{10.2pt}{40}{ਦੁਕਾਨ ਵਿਚ ਝਾਤੀ।}
\VUsaut{2.8mm}

%%K t16-11-1 p
11. --- ਖਿਡੌਣਿਆਂ ਦਾ ਸ਼ੋਅਬਾ ਛੱਡ ਕੇ ਅਸੀਂ ਕਾਠੀ-ਸਾਜ਼ੀ ਵੱਲ ਗਏ, ਤਾਂ ਜੋ \VUgras{ਕਾਠੀਆਂ}
\textsuperscript{(53)}, \VUgras{ਰਕਾਬਾਂ} \textsuperscript{(54)}, \VUgras{ਲਗਾਮਾਂ}
\textsuperscript{(55)}, \VUgras{ਦਹਾਣਿਆਂ} \textsuperscript{(56)} ਤੇ
\VUgras{ਸਵਾਰੀ ਦੇ ਚਾਬੁਕਾਂ} ਉੱਤੇ ਇਕ ਨਜ਼ਰ ਮਾਰ ਲਈਏ। ਜਦੋਂ ਤਕ
\VUgras{ਸ਼ੋਅਬੇ ਦੇ ਪ੍ਰਬੰਧਕ} \textsuperscript{(57)} ਮੇਰੇ ਦਾਦਾ ਨੂੰ ਕੁਝ ਜਾਣਕਾਰੀ ਦਿੰਦੇ
ਰਹੇ, ਮੈਂ \VUgras{ਚੀਨੀ ਮਿੱਟੀ} ਤੇ \VUgras{ਸ਼ੀਸ਼ੇ ਦੇ ਭਾਂਡਿਆਂ} \textsuperscript{(58)}
ਦੀ ਸਜਾਵਟ ਵੇਖਣ ਗਿਆ, ਜਿੱਥੇ ਸੋਹਣੇ \VUgras{ਸਲਾਦ ਦੇ ਕਟੋਰੇ} ਤੇ ਨਾਜ਼ੁਕ \VUgras{ਚਾਹ-ਦਾਨ}
\textsuperscript{(59)} ਸਨ।

%%K t16-12-1 p
12. --- ਪਹਿਲੀ ਮੰਜ਼ਿਲ ਉੱਤੇ ਚੜ੍ਹਨ ਲਈ ਅਸੀਂ \VUgras{ਕਮਰਬੰਦਾਂ} \textsuperscript{(60)}
ਦੀ ਖਿੜਕੀ ਦੇ ਪਿੱਛਿਓਂ ਲੰਘੇ, ਜੋ ਜੁਰਾਬਾਂ ਦੇ ਸ਼ੋਅਬੇ ਕੋਲ ਹੈ, ਜਿੱਥੇ ਬਹੁਤ ਸੋਹਣੀਆਂ
\VUgras{ਨਿੱਕਰਾਂ} \textsuperscript{(63)} ਵਿਛੀਆਂ ਸਨ। ਕੋਲ ਹੀ ਇਕ
\VUgras{ਦੁਕਾਨ ਦੀ ਸਹਾਇਕ} \textsuperscript{(61)} ਇਕ \VUgras{ਗਾਹਕ}
\textsuperscript{(62)} ਨੂੰ ਸੋਹਣੇ ਰੇਸ਼ਮੀ \VUgras{ਕੱਪੜੇ} \textsuperscript{(64)} ਚੁਣਨ
ਵਿਚ ਮਦਦ ਕਰ ਰਹੀ ਸੀ।

%%K t16-12-2 p
ਪੌੜੀਆਂ ਚੜ੍ਹਦੇ ਹੋਏ ਮੈਂ ਵੇਖਿਆ ਕਿ ਦੁਕਾਨ ਜਾਪਾਨੀ \VUgras{ਲਾਲਟੈਣਾਂ}
\textsuperscript{(65)}, \VUgras{ਛਤਰੀਆਂ} \textsuperscript{(66)} ਤੇ ਵੱਡੀਆਂ
\VUgras{ਖਜੂਰਾਂ} \textsuperscript{(70)} ਨਾਲ ਸਜੀ ਹੈ।

%%K t16-13-1 p
13. --- ਪਹਿਲੀ ਮੰਜ਼ਿਲ ਉੱਤੇ ਵੇਖਣ ਨੂੰ ਵੱਧ ਕੁਝ ਨਾ ਸੀ: ਸਿਰਫ਼ ਸਾਜ਼-ਸਾਮਾਨ (*),
\VUgras{ਤਿਜੌਰੀਆਂ} \textsuperscript{(67)}, \VUgras{ਦਰਾਜ਼ਾਂ ਵਾਲੀਆਂ ਅਲਮਾਰੀਆਂ}
\textsuperscript{(68)} ਤੇ \VUgras{ਬਾਲ-ਗੱਡੀਆਂ} \textsuperscript{(69)}। ਪਰ ਦੁਕਾਨ ਦਾ
ਸਾਰਾ ਦ੍ਰਿਸ਼ ਬਹੁਤ \VUgras{ਪਰਚਾਵਾ ਦੇਣ ਵਾਲਾ} ਸੀ।

%%K t16-14-1 p
14. --- ਆਪਣੇ ਹੇਠਾਂ ਮੈਨੂੰ ਸਾਜ਼ ਵਿਖਾਈ ਦਿੱਤੇ: \VUgras{ਬੀਣਾਂ}
\textsuperscript{(71)}, \VUgras{ਗਿਟਾਰ} \textsuperscript{(72)},
\VUgras{ਮੈਂਡੋਲਿਨ} \textsuperscript{(73)}, \VUgras{ਕਾਰਨੇਟ} \textsuperscript{(74)},
\VUgras{ਕਲੈਰਿਨੇਟ} \textsuperscript{(75)}, ਆਪਣੇ \VUgras{ਗਜ਼} \textsuperscript{(76)}
ਸਮੇਤ ਵਾਇਲਨ ਤੇ ਇਕ \VUgras{ਵੱਡਾ ਢੋਲ} \textsuperscript{(77)} ਤਕ। ਕੁਝ ਹੋਰ ਅੱਗੇ, ਲਿਖਣ ਦੇ
ਸਾਮਾਨ ਦੇ ਕਾਊਂਟਰ ਉੱਤੇ, ਇਕ \VUgras{ਵਿਦਿਆਰਥੀ} \textsuperscript{(78)}
\VUgras{ਦੁਕਾਨ ਦੇ ਸਹਾਇਕ} \textsuperscript{(79)} ਨਾਲ ਇਕ ਕਾਗਜ਼ ਦੇ ਥੈਲੇ ਦਾ ਮੁੱਲ-ਭਾਅ ਕਰ
ਰਿਹਾ ਸੀ। ਉਨ੍ਹਾਂ ਕੋਲ ਮੈਨੂੰ ਇਕ ਸੋਹਣੀ \VUgras{ਕਾਇਦੇ ਦੀ ਕਿਤਾਬ} \textsuperscript{(80)}
ਵਿਖਾਈ ਦਿੱਤੀ।

%%K t16-14-2 p
ਦੁਕਾਨ ਦੇ ਵਿਚਕਾਰ ਇਕ ਉੱਚਾ \VUgras{ਕ੍ਰਿਸਮਸ ਦਾ ਰੁੱਖ} \textsuperscript{(81)} ਖੜਾ ਸੀ,
ਜਿਸ ਉੱਤੇ ਚੁਫੇਰੇ ਮੋਮਬੱਤੀਆਂ, ਗਹਿਣੇ ਤੇ ਹਰ ਤਰ੍ਹਾਂ ਦੇ ਖਿਡੌਣੇ ਟੰਗੇ ਸਨ:
\VUgras{ਲੱਕੜ ਦੇ ਘੋੜੇ} \textsuperscript{(82)}, \VUgras{ਗੁੱਡੀਆਂ}
\textsuperscript{(83)}, ਵਗ਼ੈਰਾ।

%%K t16-14-3 p
\VUgras{ਅਤਰ ਦੀ ਦੁਕਾਨ} \textsuperscript{(84)}, ਆਪਣੇ \VUgras{ਦੰਦਾਂ ਦੇ ਮੰਜਨਾਂ} ਤੇ
ਆਪਣੇ ਵੱਖੋ-ਵੱਖ \VUgras{ਅਤਰਾਂ} ਸਮੇਤ, ਇਕ ਬਹੁਤ ਸੋਹਣੀ ਵਾਸ਼ਨਾ ਛੱਡਦੀ ਸੀ ਜੋ ਸਾਡੇ ਤਕ ਉੱਠ
ਆਉਂਦੀ ਸੀ।

%%K t16-tit-5 sub
\VUsaut{3.4mm}
\VUpk{10.2pt}{40}{ਅਸੀਂ ਕੁਝ ਖ਼ਰੀਦਦੇ ਹਾਂ।}
\VUsaut{2.8mm}

%%K t16-15-1 p
15. --- ਕਿਉਂਕਿ ਮੇਰੇ ਦਾਦਾ ਨੂੰ ਵੇਲਾ ਲੰਮਾ ਲੱਗਣ ਲੱਗ ਪਿਆ ਸੀ, ਅਸੀਂ ਵੱਡੀ ਪੌੜੀ ਤੋਂ ਮੁੜ
ਹੇਠਾਂ ਉਤਰੇ। ਉਹ ਇਕ \VUgras{ਤਾਰਾ-ਵਿੱਦਿਆ ਦੀ ਸੂਚੀ} \textsuperscript{(85)} ਦੇ ਸਾਹਮਣੇ ਪਲ
ਭਰ ਰੁਕੇ, ਜਿਸ ਉੱਤੇ, ਉਨ੍ਹਾਂ ਦੱਸਿਆ, \VUgras{ਪੂਛਲ ਤਾਰੇ} \textsuperscript{(a)},
\VUgras{ਚੰਨ ਤੇ ਸੂਰਜ ਦੇ ਗ੍ਰਹਿਣ} \textsuperscript{(b)}, \VUgras{ਚੌਹਾਂ ਸੇਧਾਂ}
\textsuperscript{(c)} \VUgras{(ਉੱਤਰ, ਦੱਖਣ, ਪੂਰਬ ਤੇ ਪੱਛਮ)} ਵਾਲਾ ਇਕ ਸੇਧ-ਚੱਕਰ, ਦੋਹਾਂ
\VUgras{ਅੱਧ-ਗੋਲਿਆਂ} \textsuperscript{(d)} ਦਾ ਨਕਸ਼ਾ, \VUgras{ਗ੍ਰਹਿ}
\textsuperscript{(e)} ਤੇ \VUgras{ਸੂਰਜ-ਮੰਡਲ} \textsuperscript{(f)} ਵਿਖਾਏ ਗਏ ਸਨ। ਮੈਨੂੰ
ਉਸ ਵਿਚੋਂ ਕੁਝ ਸਮਝ ਨਾ ਆਇਆ, ਤੇ ਮੈਨੂੰ ਕਿਤੇ ਵੱਧ ਦਿਲਚਸਪੀ ਉਸ
\VUgras{ਸਾਮਾਨ ਪਹੁੰਚਾਉਣ ਵਾਲੇ} \textsuperscript{(86)} ਦੀ ਗੋਟੇ ਵਾਲੀ ਵਰਦੀ ਵਿਚ ਹੋਈ ਜੋ
ਲੰਘ ਰਿਹਾ ਸੀ, ਤੇ ਆਪਣੇ ਕੋਲ ਰੱਖੀਆਂ ਉਨ੍ਹਾਂ ਸੋਹਣੀਆਂ \VUgras{ਪੱਖੀਆਂ}
\textsuperscript{(87)} ਵਿਚ, \VUgras{ਬਟੂਇਆਂ} \textsuperscript{(88)} ਤੇ
\VUgras{ਜੇਬੀ ਬਟੂਇਆਂ} \textsuperscript{(89)} ਕੋਲ।

%%K t16-16-1 p
16. --- ਮੈਂ ਕਾਊਂਟਰ ਉੱਤੇ ਕੁਝ \VUgras{ਖੰਭ} \textsuperscript{(90)} ਤੇ
\VUgras{ਪੇਟੀ ਦੇ ਬਕਸੂਏ} \textsuperscript{(91)} ਵੀ ਤੱਕੇ, ਤੇ ਉਹ ਸੋਹਣੇ
\VUgras{ਬਣਾਉਟੀ ਫੁੱਲ} \textsuperscript{(94)} ਜੋ ਟੋਕਰੀਆਂ ਵਿਚ ਸਲੀਕੇ ਨਾਲ ਸਜੇ ਸਨ।
\VUgras{ਬਨਫ਼ਸ਼ੇ}, \VUgras{ਗੁਲ-ਸ਼ਬੋ}, \VUgras{ਸੁੰਬਲ} \textsuperscript{(95)},
\VUgras{ਜੇਰੇਨੀਅਮ} \textsuperscript{(96)}, \VUgras{ਕੰਵਲ} \textsuperscript{(97)} ਤੇ
\VUgras{ਨਸਤਰਨ} \textsuperscript{(98)} ਦੀ ਨਕਲ ਬਹੁਤ ਚੰਗੀ ਸੀ।

%%K t16-tit-6 sub
\VUsaut{4.4mm}
\VUpk{10.2pt}{40}{ਦਾਦੇ ਦੀ ਸੌਗਾਤ।}
\VUsaut{2.8mm}

%%K t16-17-1 p
17. --- ਮੈਂ ਆਪਣੇ ਦਾਦਾ ਨੂੰ ਆਖਿਆ ਕਿ ਮੈਨੂੰ ਉਨ੍ਹਾਂ \VUgras{ਦੰਦਾਂ ਦੇ ਬੁਰਸ਼ਾਂ}
\textsuperscript{(92)} ਵਿਚੋਂ ਇਕ ਖ਼ਰੀਦ ਦਿਓ ਜੋ \VUgras{ਵਾਲਾਂ ਦੀਆਂ ਪਿੰਨਾਂ}
\textsuperscript{(93)} ਕੋਲ ਇਕ ਡੱਬੇ ਵਿਚ ਸਨ। ਉਹ ਮੰਨ ਗਏ, ਤੇ ਮੈਂ ਆਪ ਗੱਲੇ ਉੱਤੇ ਆਪਣੀ
ਖ਼ਰੀਦ ਦਾ ਮੁੱਲ ਤਾਰਨ ਗਿਆ, ਜਿਸ ਕੋਲ ਇਕ ਵੱਡਾ \VUgras{ਪੈਕਟ} \textsuperscript{(100)}
ਬੰਨ੍ਹਿਆ ਜਾ ਰਿਹਾ ਸੀ, ਇਕ \VUgras{ਗਾਹਕ} \textsuperscript{(99)} ਲਈ।

%%K t16-18-1 p
18. --- ਅਚਾਨਕ ਕਲਾਤਮਕ \VUgras{ਕਾਂਸੀ ਦੇ ਬੁੱਤਾਂ} \textsuperscript{(101)} ਕੋਲ ਖੜੇ ਲੋਕਾਂ
ਵਿਚ ਹਲਚਲ ਮਚ ਗਈ। ਇਕ \VUgras{ਚੋਰ} \textsuperscript{(102)} ਨੂੰ ਹੁਣੇ ਹੁਣੇ ਇਕ
\VUgras{ਦੁਕਾਨ ਦੇ ਨਿਗਰਾਨ} \textsuperscript{(103)} ਨੇ ਕਿਸੇ ਦੀ ਜੇਬ ਕੱਟਦੇ ਰੰਗੇ ਹੱਥੀਂ ਫੜ
ਲਿਆ ਸੀ। ਉਸ ਨੂੰ \VUgras{ਗ੍ਰਿਫ਼ਤਾਰ} ਕਰਨ ਨੂੰ ਸਿਪਾਹੀ ਸੱਦਿਆ ਗਿਆ, ਤੇ ਅਸੀਂ ਨਿਕਲ ਹੀ ਰਹੇ ਸੀ
ਕਿ ਮੁਜਰਮ ਨੂੰ ਥਾਣੇ ਲੈ ਜਾਇਆ ਜਾ ਰਿਹਾ ਸੀ।
\VUsaut{6.0mm}
\VUfilet{22mm}
